\documentclass{generic_report}
% Optional: \documentclass[twocolumn]{generic_report}
% The front matter stays full width in either mode.

\usepackage{amsmath,amssymb}
\usepackage{booktabs,tabularx,array,adjustbox,longtable}
\usepackage[table]{xcolor}
\usepackage{textcomp}
\usepackage{tikz}
\usetikzlibrary{positioning,arrows.meta,fit,backgrounds}
\usepackage{subcaption}
\usepackage{paralist}
\usepackage{enumitem}
\usepackage{listings}
\newif\ifusealternatefonts
\usealternatefontsfalse % Optional TeX-provided sans font; no bundled fonts.
\ifusealternatefonts\usepackage[scaled=.95]{helvet}\fi
% Provider logos appear only beside model names, never in score cells.
\newcommand{\ModelLogo}[1]{\raisebox{-1.5pt}{\includegraphics[height=1.05em]{figures/logos/models/#1.png}}\hspace{0.28em}}
\newcommand{\OpenAIModel}[1]{\ModelLogo{openai}#1}
\newcommand{\AnthropicModel}[1]{\ModelLogo{anthropic}#1}
\newcommand{\GeminiModel}[1]{\ModelLogo{gemini}#1}
\newcommand{\GrokModel}[1]{\ModelLogo{grok}#1}
\newcommand{\QwenModel}[1]{\ModelLogo{qwen}#1}
\newcommand{\GLMModel}[1]{\ModelLogo{glm}#1}
\newcommand{\KimiModel}[1]{\ModelLogo{kimi}#1}
\newcommand{\MiMoModel}[1]{\ModelLogo{mimo}#1}
\newcommand{\MiniMaxModel}[1]{\ModelLogo{minimax}#1}
\newcommand{\DeepSeekModel}[1]{\ModelLogo{deepseek}#1}
\newcommand{\ModelRef}[1]{\textsuperscript{\citealp{#1}}}

\newcolumntype{Y}{>{\raggedright\arraybackslash}X}
\setlist[itemize]{leftmargin=*,topsep=3pt,itemsep=2pt}
\setlist[enumerate]{leftmargin=*,topsep=3pt,itemsep=2pt}
\lstset{basicstyle=\ttfamily\small,breaklines=true,columns=fullflexible,
  keepspaces=true,showstringspaces=false,frame=single,rulecolor=\color{ReportThemeLight}}

\newcommand{\homepage}{\raisebox{-1.5pt}{\includegraphics[height=1em]{figures/logos/page-logo.pdf}}}
\newcommand{\github}{\raisebox{-1.5pt}{\includegraphics[height=1em]{figures/logos/github-logo.pdf}}}
\newcommand{\huggingface}{\raisebox{-1.5pt}{\includegraphics[height=1em]{figures/logos/hf-logo.pdf}}}
\newcommand{\CoreTerm}{[Core Capability]}
\newif\ifshowcontents
\showcontentstrue % Set \showcontentsfalse for a short paper.
\newif\ifshowgallery
\showgallerytrue % Template-only macro gallery; remove from a real manuscript.

\title{\CoreTerm: [Problem, Approach, and Scope]}
\teaser{\resizebox{\linewidth}{!}{%
\begin{tikzpicture}[font=\sffamily,>=Stealth]
\tikzset{
 panel/.style={draw=ReportThemeLight,fill=ReportPanel,rounded corners=5pt,minimum width=4cm,minimum height=4.2cm},
 tag/.style={draw=ReportTheme,fill=white,rounded corners=3pt,text width=3.25cm,align=center,minimum height=0.65cm,font=\small},
 title/.style={font=\bfseries\small,text=ReportTheme}
}
\node[panel] at (0,0) {};
\node[panel] at (4.65,0) {};
\node[panel] at (9.3,0) {};
\node[title] at (0,1.6) {Problem setting};
\node[title] at (4.65,1.6) {Research challenges};
\node[title] at (9.3,1.6) {Approach and evidence};
\node[tag] (q) at (0,.65) {Input: [context and task]};
\node[tag] (y) at (0,-.65) {Output: [response or plan]};
\draw[->,color=ReportTheme] (q) -- (y);
\node[tag] at (4.65,.85) {[Challenge A]};
\node[tag] at (4.65,-.1) {[Challenge B]};
\node[tag] at (4.65,-1.05) {[Challenge C]};
\node[tag] (m) at (9.3,.65) {[Method / resource]};
\node[tag] (e) at (9.3,-.65) {[Evaluation and feedback]};
\draw[->,color=ReportTheme] (m) -- (e);
\draw[->,thick,color=ReportTheme] (2.05,0) -- (2.6,0);
\draw[->,thick,color=ReportTheme] (6.7,0) -- (7.25,0);
\node[draw=ReportTheme,fill=ReportTheme,text=white,rounded corners=4pt,text width=12.8cm,align=center,minimum height=.9cm,font=\small] at (4.65,-3) {Key finding: [one verified observation and its bounded implication]};
\end{tikzpicture}%
}
}
\teasercaption{\textbf{Figure 1. Research overview.} [Paper name] studies \CoreTerm{} through [challenge A], [challenge B], and [challenge C]. Replace the schematic with the real task, approach, and one evidence-backed finding.}
\resourcelinks{\begin{center}
  \begin{tabular}{@{}l@{}}
    \homepage\ \textbf{Page}\quad \url{https://example.org/project}\\[1mm]
    \github\ \textbf{Code}\quad \url{https://example.org/code}\\[1mm]
    \huggingface\ \textbf{Data}\quad \url{https://example.org/data}
  \end{tabular}
\end{center}}

\author[1,*]{First Author}
\author[2,*]{Second Author}
\affiliation[1]{First Institution}
\affiliation[2]{Second Institution}
\contribution[*]{Equal contribution.}
\abstract{[Background and specific capability gap]. We introduce [method, dataset, or system] to study \CoreTerm{} under [setting and scope]. The approach connects [design A], [design B], and [design C] with [evaluation protocol]. Experiments on [verified evaluation setting] reveal [supported finding], with [measured benefit or diagnosed limitation]. These results provide evidence for [bounded implication]. This paragraph is an abstract-writing scaffold; every bracketed field must be replaced by verified content, and citations belong in the Introduction.}
\date{\today}
\correspondence{\email{author@example.org}}
\checkdata[Version]{[Manuscript version]}

\begin{document}
\refstepcounter{figure}\label{fig:teaser}
\maketitle
\ifshowcontents
  \tableofcontents
\fi
\ReportBeginMain

% Four supported story routes: benchmark/data, method, system/agent, and
% training/data-loop. Keep only the sections needed by the actual paper.
% The planning matrix and final review checklist are in README.md.
\section{Introduction}
\label{sec:introduction}

% P1: background -> capability gap -> importance.
[The field aims to achieve a long-term goal]. Current large language models (LLMs) can [established capability], while \CoreTerm{} requires [bounded additional capability] in [specific setting]. This distinction matters because [task-grounded consequence]. The overview in \cref{fig:teaser} connects this research question to the proposed design and its eventual evidence.

% P2: classify prior work; put citations beside the supported objects.
Existing work can be organized by [task resources], [model design], and [evaluation]. For example, the Transformer~\citep{vaswani2017attention} is a verified citation example for an architecture name. Replace this syntax demonstration with sources relevant to the actual paper and explain which objective each research direction serves.

% P3: parallel challenges, then an integrating sentence.
Three challenges define the study. First, \textbf{[Challenge A]} concerns [capability and source of difficulty]. Second, \textbf{[Challenge B]} concerns [capability and failure mechanism]. Third, \textbf{[Challenge C]} concerns [capability and evaluation consequence]. Together, \CoreTerm{} requires a design that connects these capabilities through a common task and evaluation protocol.

% P4: same order as the challenge paragraph.
In this paper, we introduce [approach] for \CoreTerm{} under [scope]. [Component A] addresses [Challenge A] by [mechanism]. [Component B] addresses [Challenge B] through [constraint]. [Component C] addresses [Challenge C] with [verification or feedback]. The formulation, construction, and evaluation are developed in \cref{sec:method}.

% P5: verified observations first, interpretation second; no invented numbers.
Experiments on [setting] show [verified quantitative observation] and [verified qualitative observation]. These findings suggest [bounded capability insight], supported by the main results, ablations, and diagnostic cases in \cref{sec:experiments}. Populate this paragraph only after the evidence is available.

% P6: compactitem; do not put an unescaped [label] immediately after \item.
Our contributions are organized around the research question:
\begin{compactitem}
\item \textbf{Task formulation.} [Define the capability, inputs, outputs, and scope.]
\item \textbf{Approach or resource.} [State the method or dataset contribution at a research level.]
\item \textbf{Evaluation and evidence.} [State the evaluation design and supported finding.]
\end{compactitem}

\FloatBarrier
\section{Related Work}
\label{sec:related-work}

\subsection{Tasks, Resources, and Capability Coverage}
[Summarize verified resources by the capability they measure]. Explain their inputs, supervision, and evaluation scope before making a comparison. \Cref{tab:coverage} illustrates a capability table whose columns should match the challenges in the Introduction.

[Explain the difference in task goal, distribution, or supervision]. The symbols below are visual placeholders, rather than evidence about real datasets. Replace every row and symbol after checking the underlying sources; place the proposed resource last only when its coverage is established.

\begin{widetable}[htbp]
\centering
\caption{\textbf{Capability comparison scaffold.} \cmark: full support; \pmark: partial support under [explicit restriction]; \xmark: no support. All entries are placeholders.}
\label{tab:coverage}
\small
\begin{tabularx}{\textwidth}{@{}Ycccc@{}}
\toprule
Resource & [Challenge A] & [Challenge B] & [Challenge C] & Train split\\
\midrule
Prior resource A & \pmark & \xmark & \cmark & \xmark\\
Prior resource B & \cmark & \pmark & \xmark & \cmark\\
\rowcolor{CustomRowShade}
\textbf{Proposed resource [replace]} & \cmark & \cmark & \cmark & \cmark\\
\bottomrule
\end{tabularx}
\end{widetable}

\subsection{Models and System Design}
Model reports provide examples of source-linked names: GPT-4~\citep{openai2023gpt4}, Gemini 1.5~\citep{gemini2024gemini15}, and Qwen2.5~\citep{qwen2024qwen25}. These historical reports are included to exercise the citation and logo macros. They make no claim about the models' performance on the placeholder task.

[Describe the relevant modeling distinction and how it changes the present task]. Keep comparisons about research objectives, capabilities, assumptions, or evidence. Replace this paragraph with a synthesis tied to the contribution rather than a chronological list of papers.

\subsection{Evaluation and Learning Feedback}
[Group evaluation and learning work into two or three relevant lines]. The textual citation form is also available: \citet{vaswani2017attention} provide the source used by this citation-format example. Select the appropriate formal publication when a paper has both a preprint and a proceedings version.

[Explain the remaining gap in verification, diagnosis, or feedback]. The intended conclusion is a precise capability boundary that can be evaluated in the Method and Experiments, with each external claim supported beside its source.

\FloatBarrier
\section{Method}
\label{sec:method}

The approach connects [construction], [decision making], and [verification] to study \CoreTerm{}. Each component addresses one of the challenges from \cref{sec:introduction}. \Cref{fig:pipeline} separates data construction from model execution and closes the analysis or learning feedback loop.

\begin{widefigure}[htbp]
\centering
\resizebox{0.94\textwidth}{!}{%
\begin{tikzpicture}[font=\small\sffamily,>=Stealth]
\tikzset{stage/.style={draw=ReportTheme,fill=ReportPanel,rounded corners=4pt,text width=3cm,minimum height=.9cm,align=center}}
\node[stage] (source) at (0,0) {Source records\\and task inputs};
\node[stage] (gen) at (4,0) {Instance generation\\and difficulty control};
\node[stage] (valid) at (8,0) {Validation, filtering,\\and leakage checks};
\node[stage] (model) at (8,-2) {Model inference\\or training};
\node[stage] (eval) at (4,-2) {Programmatic evaluation\\or task reward};
\node[stage] (feedback) at (0,-2) {Error analysis\\and repair feedback};
\node[stage,fill=white] (out) at (4,-3.8) {Verified output\\and evidence record};
\foreach \a/\b in {source/gen,gen/valid,valid/model,model/eval,eval/feedback,eval/out}
  \draw[->,thick,color=ReportTheme] (\a) -- (\b);
\draw[->,dashed,thick,color=ReportTheme] (feedback.west) -- ++(-.7,0) |- (gen.north);
\end{tikzpicture}%
}

\caption{\textbf{End-to-end pipeline scaffold.} Source records become validated instances, then pass through inference or training and evaluation. Feedback returns to repair or learning; replace these labels and edges with the real scientific design.}
\label{fig:pipeline}
\end{widefigure}

\subsection{Task Formulation}
Let $q$ denote the input, $c$ the task context, $y^\star$ the verified target, and $\mathcal{C}$ the allowed constraints. An instance is $d=(q,c,y^\star,\mathcal{C})$, and the dataset is $\mathcal{D}=\{d_i\}_{i=1}^{N}$. A model returns reasoning $r$ and a final response $\hat y$:
\begin{equation}
\begin{aligned}
(r,\hat y) &= f_\theta(q,c),\\
v &= \operatorname{Valid}(\hat y,\mathcal{C}),\\
s &= v\cdot g(\hat y,y^\star).
\end{aligned}
\label{eq:formulation}
\end{equation}
Here $v\in\{0,1\}$ indicates format and constraint validity, while $g$ is a task-specific scoring function. Explain why these variables capture the task, define the score range, and use the same symbols in the evaluator and analyses. \Cref{eq:formulation} is a formulation scaffold, not a prescribed evaluation metric.

\subsection{Source Construction and Instance Generation}
To address [Challenge A], construct instances from [documented sources] using [schema and sampling rule]. Specify the transformation, preserve source identifiers, and enforce [construction constraint]. If LLMs generate candidates, state the prompt/schema, automatic checks, model or human review, and repair/discard rules.

\subsection{Difficulty and Quality Control}
Describe difficulty through [context length], [dependencies], [parameter perturbations], [negative examples], or [compositional split]. Quality control should document format checks, syntax parsing, deduplication, leakage detection, annotation consistency, pilot testing, human audits, repair decisions, and dataset versioning as applicable.

\begin{table}[htbp]
\centering
\caption{\textbf{Data distribution scaffold.} Counts, lengths, difficulty, and attributes must explain how the split supports the capability claim.}
\label{tab:distribution}
\small
\begin{tabularx}{\linewidth}{@{}lccYY@{}}
\toprule
Split & Count & Avg.\ length & Difficulty & Attribute\\
\midrule
Train & [N] & [L] & [level] & [coverage]\\
Dev & [N] & [L] & [level] & [coverage]\\
Test & [N] & [L] & [level] & [held-out factor]\\
\bottomrule
\end{tabularx}
\end{table}

\Cref{tab:distribution} should be followed by an explanation of category balance, train/test relationships, deduplication, leakage prevention, and difficulty coverage. Report all counts from the actual processed version, with evidence for every filtering decision.

\subsection{Core Component and Constraints}
To address [Challenge B], use [component] to maintain [state or dependency]. The mechanism in \cref{fig:component} illustrates a technical core rather than repeating the pipeline. Define its inputs, state update, guard, and failure behavior, then explain the role of each constraint.

\begin{widefigure}[htbp]
\centering
\resizebox{0.94\textwidth}{!}{%
\begin{tikzpicture}[font=\small\sffamily,>=Stealth]
\tikzset{module/.style={draw=ReportTheme,fill=ReportPanel,rounded corners=3pt,text width=2.2cm,minimum height=.9cm,align=center}}
\node[module] (candidate) at (0,0) {Candidate\\{[action / answer]}};
\node[module] (update) at (3.2,0) {State update\\{[transition rule]}};
\node[module] (check) at (6.4,0) {Constraint check\\{[guard]}};
\node[module,fill=white] (accept) at (9.6,0) {Accepted output\\{[evidence]}};
\node[module,fill=white] (repair) at (3.2,-1.8) {Repair or reject\\{[bounded policy]}};
\draw[->,thick,color=ReportTheme] (candidate) -- (update);
\draw[->,thick,color=ReportTheme] (update) -- (check);
\draw[->,thick,color=ReportTheme] (check) -- node[above,font=\scriptsize] {valid} (accept);
\draw[->,dashed,color=ReportTheme] (check.south) |- node[below,pos=.7,font=\scriptsize] {invalid} (repair);
\draw[->,dashed,color=ReportTheme] (repair.west) -| (candidate.south);
\end{tikzpicture}%
}

\caption{\textbf{Core-component scaffold.} A state update is checked before acceptance; rejected candidates follow the explicit repair path. This diagram should be replaced with the actual mechanism or algorithm.}
\label{fig:component}
\end{widefigure}

\subsection{Evaluation Protocol}
For [Challenge C], state what each metric measures, why it fits the task, which errors it penalizes, and which irrelevant differences it ignores. Define invalid answers, whether they count as wrong, the parsing boundary, retries, and cache behavior. Separate final-answer extraction from reasoning text when the protocol requires it.

Document the test size, exact model versions, prompt and output formats, decoding, token limits, number of runs, averaging, uncertainty, and fair baseline settings. Full settings and parser scaffolds appear in \cref{app:settings,app:protocol}; these placeholders must become an executable, reproducible protocol for the actual paper.

\FloatBarrier
\section{Experiments}
\label{sec:experiments}

\subsection{Experimental Setup}
Specify [datasets and splits], [exact model variants], [prompt/output formats], [decoding], [token budgets], [number of runs], [statistical aggregation], [invalid handling], and [retry/cache policies]. Keep baseline settings comparable and identify every intentional difference. Training experiments additionally require the hyperparameters in \cref{tab:hyperparameters}.

\subsection{Main Results}
\Cref{tab:main-results} demonstrates a compact result table with provider logos, source citations, directional text arrows, restrained row grouping, color cells, and best/second formatting. Every value is deliberately marked \texttt{TBD}; the shade intensities and rank marks only preview typography and must be recomputed from real measurements.

\begin{widetable}[htbp]
\centering
\caption{\textbf{Main-result layout scaffold.} Task scores use higher-is-better arrows; invalid rate uses lower-is-better. All values, color intensities, and ranking marks are placeholders. Historical model names are citation examples, with no reported comparison.}
\label{tab:main-results}
\begingroup
\small
\setlength{\tabcolsep}{6pt}
\renewcommand{\arraystretch}{1.12}
\begin{adjustbox}{max width=\textwidth}
\begin{tabular}{@{}lcccc@{}}
\toprule
Model / system & Overall \MetricUp & Task A \MetricUp & Task B \MetricUp & Invalid \MetricDown\\
\midrule
\rowcolor{CustomRowShade}
\multicolumn{5}{l}{\textbf{Hosted model examples}}\\
\OpenAIModel{GPT-4}\ModelRef{openai2023gpt4} & \ScoreCell{25}{TBD} & \SecondScore{40}{TBD} & \ScoreCell{25}{TBD} & TBD\\
\GeminiModel{Gemini 1.5}\ModelRef{gemini2024gemini15} & \SecondScore{40}{TBD} & \ScoreCell{25}{TBD} & \SecondScore{40}{TBD} & TBD\\
\rowcolor{CustomRowShade}
\multicolumn{5}{l}{\textbf{Open-weight and proposed system examples}}\\
\QwenModel{Qwen2.5}\ModelRef{qwen2024qwen25} & \ScoreCell{25}{TBD} & \ScoreCell{25}{TBD} & \ScoreCell{25}{TBD} & TBD\\
\textbf{Proposed system [replace]} & \BestScore{55}{TBD} & \BestScore{55}{TBD} & \BestScore{55}{TBD} & TBD\\
\bottomrule
\end{tabular}
\end{adjustbox}
\endgroup
\end{widetable}

The results paragraph should report [measured range], [comparison], [uncertainty], and [invalid-output frequency] from \cref{tab:main-results}. State the observation before interpreting it; use one color scale for comparable columns and avoid implying that two unrelated metrics have a common scale.

\paragraph{Insight 1: [Evidence-backed capability finding].}
Explain why the measured observation supports \CoreTerm{} within [scope]. Distinguish the relevant capability from simpler alternative explanations using a controlled comparison, and identify any remaining limitation.

\subsection{Ablation Studies}
\Cref{tab:ablations} illustrates two complementary ablations. Removing a component must test its corresponding design claim; changing validation or feedback must isolate correctness or robustness. Keep other settings fixed, describe the expected diagnostic relationship, and interpret only measured changes.

\begin{widetable}[htbp]
\centering
\caption{\textbf{Ablation scaffolds.} Replace placeholders with controlled results and map each removed mechanism to a challenge.}
\label{tab:ablations}
\begin{subtable}[t]{.48\textwidth}
\centering\small
\caption{Component contribution.}
\label{tab:component-ablation}
\begin{tabular}{@{}lcc@{}}
\toprule
Variant & Overall \MetricUp & [Capability] \MetricUp\\
\midrule
Without component A & \DimCell{15}{TBD} & TBD\\
Without component B & \DimCell{15}{TBD} & TBD\\
Full approach & \BestScore{35}{TBD} & TBD\\
\bottomrule
\end{tabular}
\end{subtable}\hfill
\begin{subtable}[t]{.48\textwidth}
\centering\small
\caption{Validation and feedback.}
\label{tab:validation-ablation}
\begin{tabular}{@{}lcc@{}}
\toprule
Variant & Valid \MetricUp & Robust \MetricUp\\
\midrule
Without validation & TBD & TBD\\
Without feedback & TBD & TBD\\
Full approach & TBD & TBD\\
\bottomrule
\end{tabular}
\end{subtable}
\end{widetable}

Describe the changes in \cref{tab:component-ablation,tab:validation-ablation}, then explain which mechanism the comparison supports and which confounders remain. Avoid treating an ablation label alone as causal evidence.

\subsection{Diagnostic Analysis}
\Cref{fig:analysis} supplies three editable analysis layouts: error taxonomy, difficulty with uncertainty, and a budget/scaling trend. Their geometry is schematic, with no numeric observations. Replace them with measured arrays and preserve the corresponding labels, units, uncertainty definition, and data provenance.

\begin{widefigure}[htbp]
\centering
\begin{subfigure}[t]{.32\textwidth}
\centering
\resizebox{\linewidth}{!}{%
\begin{tikzpicture}[font=\scriptsize\sffamily,>=Stealth]
\draw[->] (0,0) -- (3.8,0) node[right] {type};
\draw[->] (0,0) -- (0,2.5) node[above] {[rate]};
\foreach \x/\name in {.45/drift,1.55/order,2.65/invalid} {
  \draw[draw=ReportTheme,fill=ReportPanel] (\x,0) rectangle ++(.65,1.35);
  \node[rotate=25,anchor=north east] at (\x+.6,-.15) {\name};
}
\node[text=ReportTheme] at (1.9,2) {replace with measured bars};
\end{tikzpicture}}
\caption{Error taxonomy scaffold.}
\end{subfigure}\hfill
\begin{subfigure}[t]{.32\textwidth}
\centering
\resizebox{\linewidth}{!}{%
\begin{tikzpicture}[font=\scriptsize\sffamily,>=Stealth]
\draw[->] (0,0) -- (3.8,0) node[right] {difficulty};
\draw[->] (0,0) -- (0,2.5) node[above] {[metric]};
\foreach \x/\name in {.6/low,1.8/mid,3/high} {
 \draw[ReportTheme] (\x,.6) -- (\x,1.8);
 \fill[ReportTheme] (\x,1.2) circle (2pt);
 \node[below] at (\x,-.1) {\name};
}
\node[text=ReportTheme] at (1.9,2.1) {replace points and intervals};
\end{tikzpicture}}
\caption{Difficulty and uncertainty.}
\end{subfigure}\hfill
\begin{subfigure}[t]{.32\textwidth}
\centering
\resizebox{\linewidth}{!}{%
\begin{tikzpicture}[font=\scriptsize\sffamily,>=Stealth]
\draw[->] (0,0) -- (3.8,0) node[right] {budget};
\draw[->] (0,0) -- (0,2.5) node[above] {[metric]};
\draw[dashed,thick,ReportTheme] (.4,1.2) -- (3.3,1.2);
\foreach \x in {.4,1.35,2.3,3.3} \draw[ReportTheme,fill=white] (\x,1.2) circle (2.5pt);
\node[text=ReportTheme] at (1.9,2.1) {replace with measured series};
\end{tikzpicture}}
\caption{Scaling or resource trend.}
\end{subfigure}

\caption{\textbf{Three analysis scaffolds.} Bars, intervals, and the dashed series are layout placeholders; no empirical trend is asserted. Replace each panel with verified data and a specific diagnostic question.}
\label{fig:analysis}
\end{widefigure}

\paragraph{Insight 2: [Named error mechanism].}
Describe the phenomenon, its likely cause, a supporting example, and its consequence for the central claim. Candidate error names include parameter drift, order inversion, missing steps, dependency breaks, and invalid format; use only categories actually observed.

\paragraph{Insight 3: [Difficulty or resource boundary].}
Explain how the measured trend changes with [difficulty, scale, or budget], state how uncertainty was estimated, and connect the boundary to the relevant design constraint. Supplementary measurements belong in \cref{app:supplement}.

\subsection{Case Study}
\Cref{fig:cases} pairs two compressed cases, with the sample ID, task context, gold, output excerpt, score, and diagnosis. The main text should explain how these cases support the measured failure mode. Complete inputs, outputs, rubrics, and score records appear in \cref{app:case-a,app:case-b}.

\begin{widefigure}[htbp]
\centering
\begin{minipage}[t]{.485\textwidth}
\begin{tcolorbox}[enhanced,title={Case A: [success behavior]},fontupper=\small,colback=LighterGray,colframe=DeepPurple,colbacktitle=DeepPurple,coltitle=White,left=3mm,right=3mm]
\textbf{ID:} [case A]\\
\textbf{Input:} [one-line context]\\
\textbf{Gold:} [verified answer]\\
\textbf{Output excerpt:} [short excerpt]\\
\textbf{Score:} [verified score]\\
\textbf{Diagnosis:} [preserved constraint]
\end{tcolorbox}
\end{minipage}\hfill
\begin{minipage}[t]{.485\textwidth}
\begin{tcolorbox}[enhanced,title={Case B: [failure mode]},fontupper=\small,colback=LighterGray,colframe=DeepPurple,colbacktitle=DeepPurple,coltitle=White,left=3mm,right=3mm]
\textbf{ID:} [case B]\\
\textbf{Input:} [one-line context]\\
\textbf{Gold:} [verified answer]\\
\textbf{Output excerpt:} [short excerpt]\\
\textbf{Score:} [verified score]\\
\textbf{Diagnosis:} [specific broken dependency]
\end{tcolorbox}
\end{minipage}

\caption{\textbf{Compressed case-study scaffold.} Use one verified success and one verified failure to expose the capability distinction. All fields are placeholders; the complete records belong in the appendix.}
\label{fig:cases}
\end{widefigure}

\FloatBarrier
\section{Conclusion}
\label{sec:conclusion}

We studied \CoreTerm{} under [task setting] through [method, dataset, or system]. The evaluation supports [verified finding] and identifies [bounded capability insight], linking the research question to the design, measurements, and cases already presented. Future work can investigate [specific unresolved direction] within the limitations of [data, protocol, or model coverage]. Replace this scaffold with a concise conclusion that introduces no new evidence.

\FloatBarrier

\bibliographystyle{plainnat}
\bibliography{main}

\clearpage
% Appendices use the full width in the optional two-column mode.
\onecolumn
\beginappendix
\section{Training and Experimental Settings}
\label{app:settings}

This appendix records the reproducibility details referenced in \cref{sec:method,sec:experiments}. Fill the settings in \cref{tab:settings,tab:hyperparameters} from the executed experiment configuration. Preserve exact model versions, data revisions, prompts, and parsing behavior rather than inferring defaults.

\begin{longtable}{@{}p{.28\textwidth}p{.66\textwidth}@{}}
\caption{\textbf{Evaluation settings scaffold.} Each field must describe the actual run.}\label{tab:settings}\\
\toprule
Field & Recorded value or rule\\
\midrule
\endfirsthead
\toprule
Field & Recorded value or rule (continued)\\
\midrule
\endhead
Test set and revision & [dataset name, revision, test count, split definition]\\
Models and versions & [exact model IDs, weights or endpoint version, access date]\\
Prompt template & [system/user text and task-specific fields]\\
Output schema & [required fields, answer boundary, permitted code blocks]\\
Decoding & [temperature, top-p, seed, stopping rule]\\
Token limits & [input budget, output budget, truncation handling]\\
Parsing & [final-answer extraction, schema checks, numeric normalization]\\
Invalid handling & [definition, count, scoring rule, retry eligibility]\\
Runs and aggregation & [run count, averaging, uncertainty, paired tests if used]\\
Retries and caches & [retry budget, cache keys, replay policy, leakage controls]\\
Baseline fairness & [shared budgets and prompts; every intentional difference]\\
Hardware and software & [accelerator, software versions, precision and batch size]\\
Data quality & [deduplication, leakage, annotation consistency, human audit]\\
Source provenance & [source IDs, filtering evidence, repair/discard criteria]\\
Release revision & [dataset/code revision and reproducibility entry points]\\
\bottomrule
\end{longtable}

\begin{table}[htbp]
\centering
\caption{\textbf{Training hyperparameter scaffold.} Retain this table only for experiments involving training or adaptation.}
\label{tab:hyperparameters}
\small
\begin{tabularx}{.88\textwidth}{@{}lY@{}}
\toprule
Hyperparameter & Value\\
\midrule
Base model and initialization & [exact model and revision]\\
Learning rate / schedule & [value and schedule]\\
Batch size / accumulation & [values]\\
Steps / epochs & [value and stopping rule]\\
Optimizer / regularization & [settings]\\
Sequence length / precision & [settings]\\
Training data and split & [revision and leakage safeguards]\\
Inference temperature / decoding & [settings]\\
\bottomrule
\end{tabularx}
\end{table}

\Needspace{24\baselineskip}
\subsection{Prompt, Parsing, and Invalid Outputs}
\label{app:protocol}

The prompt and parser below are placeholders for a documented protocol. Replace them with the exact text and executed rules, including how invalid outputs are counted, whether retries are allowed, and how cached responses are reused.

\begin{lstlisting}[caption={Prompt-format scaffold.},label={lst:prompt}]
SYSTEM: [Task scope and permitted tools.]
USER:
  Context: [input context]
  Question: [task]
  Output contract:
    {"answer": "[final answer]", "evidence": "[support]"}
\end{lstlisting}

\begin{lstlisting}[caption={Evaluation-rule scaffold; describe the actual parser.},label={lst:parser}]
1. Extract only the protocol-defined final answer.
2. Parse the required schema and check task constraints.
3. If invalid, record the error category and apply [rule].
4. If valid, compute the documented task score.
5. Record retry/cache provenance and aggregate all runs.
\end{lstlisting}

\Cref{lst:prompt,lst:parser} should match the submitted evaluation code. Document any accepted formatting variation so that superficial differences do not change a score, and isolate semantic errors from invalid-format errors.

\section{Supplementary Results}
\label{app:supplement}

Use this section for additional ablations, category or difficulty slices, uncertainty estimates, invalid-output distributions, and agreement analyses. \Cref{tab:supplement} is a supplementary table scaffold; name the question it answers and connect it to one of the main-text insights.

\begin{table}[htbp]
\centering
\caption{\textbf{Supplementary diagnostic scaffold.} Replace all placeholder values with measured statistics and explain the uncertainty definition.}
\label{tab:supplement}
\small
\begin{tabularx}{.92\textwidth}{@{}Yccc@{}}
\toprule
Slice / category & Count & Metric \MetricUp & Invalid \MetricDown\\
\midrule
{[Difficulty group A]} & [N] & [value and interval] & [rate]\\
{[Difficulty group B]} & [N] & [value and interval] & [rate]\\
{[Held-out composition]} & [N] & [value and interval] & [rate]\\
\bottomrule
\end{tabularx}
\end{table}

\FloatBarrier
\section{Full Case A: Success Record}
\label{app:case-a}

The compact case in \cref{fig:cases} must link to its full record. The following breakable box demonstrates the fields from the skill. Bracketed content is intentionally empty of scientific claims; replace every field with the complete evidence for a real case.

\begin{ReportCaseBox}{Case A: [sample identifier]}
\CaseFirstSection{Meta Info}
\begin{itemize}
\item \textbf{System / model:} [system and exact model]
\item \textbf{Run / seed:} [run identifier and seed]
\item \textbf{Total score:} [score and range]
\item \textbf{Duration / cost:} [seconds] / \$[cost]
\end{itemize}

\CaseSection{Task}
\textbf{Input:} [Full task input and context.]\par
\textbf{Gold / scientific goal:} [Verified target, acceptable alternatives, and task objective.]\par
\textbf{Expected output:} [Complete required answer or schema.]\par

\CaseSection{Data}
\begin{itemize}
\item \texttt{[file\_name]}: [data type, provenance, and full relevant content]. Relative path: \texttt{[relative/path/to/file]}.
\end{itemize}

\CaseSection{Rubrics}
\begin{enumerate}
\item \textbf{Text | Weight([w]):} [criterion]. \emph{Expected evidence:} [required content and validation rule].
\item \textbf{Image | Weight([w]):} [criterion]. \emph{Expected evidence:} [visual requirement and validation rule].
\end{enumerate}

\CaseSection{Generated Report}
\textbf{[Report title]}\par
[Paste the full output or report here. Preserve the complete input, reasoning or review evidence required by the protocol, and final answer. If long, continue in another named box or page; do not silently shorten the record.]

\CaseSection{Figures}
\begin{center}
\includegraphics[width=.88\linewidth]{\detokenize{figures/cases/example.pdf}}
\par\footnotesize Case evidence placeholder
\end{center}
\textbf{Evidence figure.} [Explain the actual image and how it supports the case; replace this generic vector placeholder.]

\CaseSection{Score Items}
\begin{enumerate}
\item \textbf{Text | Weight([w]) | Score([score]):} [criterion]. \emph{Reasoning:} [complete evidence-grounded assessment].
\item \textbf{Image | Weight([w]) | Score([score]):} [criterion]. \emph{Reasoning:} [complete evidence-grounded visual assessment].
\end{enumerate}

\CaseSection{Diagnosis}
[Explain the preserved constraint or capability and its relationship to a main-text metric or insight. Include the complete review and scoring record.]
\end{ReportCaseBox}

\FloatBarrier
\Needspace{20\baselineskip}
\section{Full Case B: Failure Record}
\label{app:case-b}

This second record contrasts with Case A using the same fields, allowing the reader to check the claimed failure mechanism. The box structure can span pages; continue the report or score items when the real record requires more space.

\begin{ReportCaseBox}{Case B: [sample identifier]}
\CaseFirstSection{Meta Info}
\begin{itemize}
\item \textbf{System / model:} [system and exact model]
\item \textbf{Run / seed:} [run identifier and seed]
\item \textbf{Total score:} [score and range]
\item \textbf{Duration / cost:} [seconds] / \$[cost]
\end{itemize}

\CaseSection{Task}
\textbf{Input:} [Full input and context.]\par
\textbf{Gold / scientific goal:} [Complete verified target and objective.]\par
\textbf{Expected output:} [Complete answer contract.]\par

\CaseSection{Data}
\begin{itemize}
\item \texttt{[file\_name]}: [type, provenance, and complete case data]. Relative path: \texttt{[relative/path/to/file]}.
\end{itemize}

\CaseSection{Rubrics}
\begin{enumerate}
\item \textbf{Criterion | Weight([w]):} [criterion and expected evidence].
\end{enumerate}

\CaseSection{Generated Report}
[Paste the entire model output and final answer. Preserve the exact portion containing the error together with its full context; do not substitute a representative excerpt for the full record.]

\CaseSection{Figures}
[Insert the complete relevant visual evidence, using at most two or three images per box. Use additional named boxes for further evidence.]

\CaseSection{Score Items}
\begin{enumerate}
\item \textbf{Criterion | Weight([w]) | Score([score]):} [criterion]. \emph{Reasoning:} [full assessment and cited case evidence].
\end{enumerate}

\CaseSection{Diagnosis}
[Name the observed mechanism, locate the incorrect output, explain the broken constraint, and connect it to the main-text analysis. Retain the full review and scoring record.]
\end{ReportCaseBox}

\ifshowgallery
\clearpage
\section{Template Formatting Gallery}
\label{app:gallery}

This template-only appendix verifies the bundled provider wrappers and typography. Remove the gallery from a real paper; it is a rendering aid and conveys no model evaluation. The ten logos in \cref{tab:logo-gallery} use the same baseline and height defined in \texttt{seed/model-logos.tex}.

\begin{table}[htbp]
\centering
\caption{\textbf{Provider-logo macro gallery.} Provider marks identify names; they do not replace source citations in a research comparison.}
\label{tab:logo-gallery}
\begin{tabular}{@{}ll@{}}
\toprule
Provider wrapper & Provider wrapper\\
\midrule
\OpenAIModel{OpenAI} & \AnthropicModel{Anthropic}\\[2pt]
\GeminiModel{Gemini} & \GrokModel{Grok}\\[2pt]
\QwenModel{Qwen} & \GLMModel{GLM}\\[2pt]
\KimiModel{Kimi} & \MiMoModel{MiMo}\\[2pt]
\MiniMaxModel{MiniMax} & \DeepSeekModel{DeepSeek}\\
\bottomrule
\end{tabular}
\end{table}

\begin{table}[htbp]
\centering
\caption{\textbf{Table-style macro gallery.} The visual strengths below demonstrate syntax only.}
\label{tab:style-gallery}
\begin{tabular}{@{}llll@{}}
\toprule
Normal cell & Best mark & Second mark & Dim cell\\
\midrule
\ScoreCell{15}{[value]} & \BestScore{35}{[value]} & \SecondScore{25}{[value]} & \DimCell{10}{[value]}\\
\bottomrule
\end{tabular}
\end{table}

\Cref{tab:style-gallery} exercises all score-cell wrappers. Attribute symbols use \cmark{} for full support, \pmark{} for partial support with a stated restriction, and \xmark{} for no support. Directional text macros render as \MetricUp{} and \MetricDown{} and remain outside math mode in scaled table headings.
\fi

\end{document}
